\documentclass{Preport}

\usepackage{hyperref}
\hypersetup{
    colorlinks=true,
}
\usepackage{float}
\usepackage{amsmath}

\exname{铁磁材料居里温度测量} %实验名称
\instructor{陈飞燕} %指导教师
\class{电科2502} %班级
\name{管云哲} %姓名
\stuid{3250104519} %学号

\nyear{2026} %年
\nmonth{9} %月
\nday{29} %日
\nweekday{二} %星期几

\begin{document}
\setcounter{page}{1}
\makecover

\section{预习报告（10分）}
\subsection{实验综述（5分）}
（自述实验现象、实验原理和实验方法，包括必要的光路图、电路图、公式等。不超过500字。）

铁磁质内部相邻原子间存在强烈的交换耦合作用，使电子自旋磁矩在小区域内自发整齐排列，形成磁畴，因而在无外磁场时即具有自发磁化。温度升高时，热运动逐渐破坏磁畴的有序排列；当温度达到居里温度 $T_C$ 时，磁畴瓦解，材料由铁磁态转变为顺磁态，磁导率 $\mu_r$ 由远大于 1 骤降至接近 1。$T_C$ 仅由化学成分与晶体结构决定，与晶粒大小、取向、应力等结构因素几乎无关，故称结构不灵敏参数。

磁介质的磁化规律为
\begin{equation}
B=\mu_0(H+M)=(\chi_m+1)\mu_0H=\mu_r\mu_0H=\mu H
\end{equation}

本实验采用交流电桥法测量居里温度。将铁氧体样品插入 $L_2$ 线圈后，线圈电感 $L=\mu N^2A/l\propto\mu(T)$ 随温度变化，破坏电桥平衡，桥臂间电位差 $U_{CD}$ 随之改变。取 $R_a=R_b$，电桥平衡条件为
\begin{equation}
R_w=\frac{R_b}{R_a}R_1,\qquad L_2=\frac{R_b}{R_a}L_1
\end{equation}

当电感变化 $\Delta L_2$ 时，桥路输出
\begin{equation}
\Delta U_{CD}\propto\frac{\partial U_{CD}}{\partial L_2}\Delta L_2\propto\frac{\omega\Delta L_2}{(R_b+R_w+j\omega L_2)^2}
\end{equation}

可见 $U_{CD}$ 非线性但单调地反映 $\mu$ 随温度的变化。缓慢升温（或降温），逐点记录温度 $T$ 与电压 $U$，作 $U\sim T$ 曲线，其斜率最大处（拐点）对应的温度即为居里温度 $T_C$。

\begin{figure}[H]
    \centering
    \includegraphics[width=0.62\linewidth]{mu_T_curve.png}
    \caption{磁导率 $\mu$ 与磁场 $H$ 及温度 $T$ 的关系（右图突变点对应居里温度 $T_C$）}
    \label{fig:muT}
\end{figure}

\subsection{实验重点（3分）}
（简述本实验的学习重点，不超过100字。）

理解磁畴与交换耦合作用下铁磁—顺磁转变的微观机理；掌握 RL 交流电桥的平衡条件与输出电压随电感变化的规律；学会由 $U\sim T$ 曲线确定居里温度并分析升温速率、信号频率对测量结果的影响。

\subsection{实验难点（2分）}
（简述本实验的实现难点，不超过100字。）

调平电桥需将电压调至 $1.0$\,mV 以下，需耐心微调 $R_w$；样品与温度传感器之间存在热滞后，快速加热会使测得 $T_C$ 偏高，须保持动态平衡；居里点附近 $U$ 变化剧烈，应加密取点（0.5\,$^\circ$C）以免漏掉拐点。

{\fangsong 注意事项：
\begin{enumerate}[label=\arabic*.]
    \item 用PDF格式上传“实验报告”，文件名：学生姓名+学号+实验名称+周次。
    \item “实验报告”必须递交在“学在浙大”本课程内对应实验项目的“作业”模块内。
    \item “实验报告”成绩必须在“浙江大学物理实验教学中心网站”-“选课系统”内查询。
    \item 教学评价必须在“浙江大学物理实验教学中心网站”-“选课系统”内进行，学生必须进行教学评价，才能看到实验报告成绩，教学评价须在本次实验结束后3天内进行。
\end{enumerate}}
\begin{center}
    \fbfangsong{浙江大学物理实验教学中心制}
\end{center}
\end{document}
